\documentclass[conference]{IEEEtran}
\IEEEoverridecommandlockouts
% The preceding line is only needed to identify funding in the first footnote. If that is unneeded, please comment it out.
\usepackage{cite}
\usepackage{amsmath,amssymb,amsfonts}
\usepackage{algorithmic}
\usepackage{graphicx}
\usepackage{textcomp}
\usepackage{xcolor}
\usepackage{indentfirst}
\usepackage{supertabular}
\usepackage{float}
\usepackage{enumitem}
\usepackage{array}
\usepackage{tabularx}

\usepackage{hyperref}
\def\BibTeX{{\rm B\kern-.05em{\sc i\kern-.025em b}\kern-.08em
    T\kern-.1667em\lower.7ex\hbox{E}\kern-.125emX}}
\begin{document}

\title{GomMold\\
{\small An AI-based Application for Mold Detection and Prevention}
}

\author{\IEEEauthorblockN{Aina Nur Insyeerah \\ Binti Abd Rahman}
\IEEEauthorblockA{\textit{dept. Information Systems} \\
\textit{Hanyang University}\\
Seoul, Republic of Korea \\
syeera@hanyang.ac.kr}
\and
\IEEEauthorblockN{Nur Ain Syafiqah \\ Binti Rahiman}
\IEEEauthorblockA{\textit{dept. Information Systems} \\
\textit{Hanyang University}\\
Seoul, Republic of Korea \\
ainrahiman@hanyang.ac.kr}
\footnotesize
\and
\IEEEauthorblockN{Nur Shaqeerah \\ Binti Shamsul Anuar}
\IEEEauthorblockA{\textit{dept. Information Systems} \\
\textit{Hanyang University}\\
Seoul, Republic of Korea \\
qeerashamsul@hanyang.ac.kr}
\and
\IEEEauthorblockN{Yumni Karmila \\ Binti Yunus}
\IEEEauthorblockA{\textit{dept. Information Systems} \\
\textit{Hanyang University}\\
Seoul, Republic of Korea \\
yumnikarmila@hanyang.ac.kr}
}

\maketitle
\begin{abstract}
Mold growth is a common but often overlooked problem in humid environments, leading to health risks, unpleasant odors, and structural damage in homes. GomMold is an application that introduces an AI-powered system to detect and prevent household mold that enables users to identify mold through photos. Traditional mold testing methods that rely on manual visual inspection are often time-consuming, subjective, and prone to error. By utilizing advancements in artificial intelligence (AI) and machine learning, this project applies computer vision and pattern recognition techniques to automate mold detection with high precision.

\newline The system integrates AI-driven image analysis and conversational feedback to provide users with accurate detection, and practical cleaning and prevention advice. In the future, there is potential for GomMold to be integrated into smart home devices such as humidifiers or air purifiers for automated humidity control and prevention advice.
\end{abstract}

\begin{IEEEkeywords}
Mold detection, indoor air quality, computer vision, deep learning, environmental monitoring, AI health technology, image classification, preventive maintenance
\end{IEEEkeywords}

\section{Role Assignment}
The distribution of project roles is summarized in Table 1 to optimize workflow within a small team.  

\begin{table}[h]
\setlength{\tabcolsep}{12pt}
\renewcommand{\arraystretch}{1.5}
\begin{tabular}{p{1.3cm} l p{3.2cm}} \hline
Roles & Name & Task description \& etc.\\ \hline
User & Nur Ain Syafiqah & 
The primary role is to experience and validate the AI mold detection model on application from a non-experts' perspective. Also, play the role to define real-world user scenarios by detecting and inspecting mold on various surfaces such as walls, furniture, and bathrooms.\\
AI Developer & Nur Shaqeerah & Responsible for artificial intelligence mold detection from developing, training and optimizing the image recognition and classification model, to accurately detect and identify the mold from user-uploaded photos.\\

\end{tabular}
\end {table}

\begin{table}[h]
\setlength{\tabcolsep}{12pt}
\renewcommand{\arraystretch}{1.5}
\begin{tabular}{p{1.3cm} l p{3.2cm}}
Software developer & Aina Nur Insyeerah & Responsible for building the technical framework that supports the AI to integrate in this project from designing the frontend, backend system and database. \\
Development manager & Yumni Karmila & Oversee and coordinate all aspects of the AI mold detection app’s development project. Additionally, responsible for timeline management, project planning and making sure the frontend, backend and AI integrate smoothly without any problems.  \\ 
\hline
\end{tabular}
\end{table} 

\section{Introduction}
\subsection{Motivation}

{\it Moving to South Korea as foreign students from Malaysia, one challenge that surprised us the most was not the food or the culture, it was the persistent spots of mold appearing in our homes. Especially during the summer season where it’s humid and raining all day long and in my case even during winter because of the condensation from the window dripped into my wallpaper, leaving it all wet. Unlike South Korea, we do not experience four seasons or drastic temperature changes in Malaysia. Additionally, house ventilation systems in South Korea also differ from those in Malaysia, which further affects how mold develops indoors. At first, we did not know whether these spots were actually mold or just dirt clumping up, as we were unfamiliar with it. It was only when the situation worsened that we realized that it was too late. If we had recognized that it was mold earlier, we could have prevented it from getting severe. We realized that we are not the only ones facing this problem. Students living in rented houses, poorly ventilated spares, and even long-term residents in Korea often struggle to identify mold growth. It is not just a lack of knowledge, many people, even those who have lived in Korea and experienced it for years, have difficulty recognizing mold in their homes.

Our motivation for this project comes from this experience. We wanted to create a tool that helps anyone to detect and manage mold confidently. We aim to develop an app that is user-friendly and allows everyone to easily identify mold in their homes within seconds. By combining AI with accessible technology, we aim to simplify a problem that is often overlooked but can have serious health and structural consequences. With this project, users can identify whether it is mold from photos taken from their phones and receive clear cleaning advice all in one place. Our goal is not just to make mold detection easier, but to empower people to live in healthier, safer homes, turning an unfamiliar and potentially hazardous issue into something manageable and understandable. We are also motivated by the potential of AI technology to make this process intuitive and user-friendly. By leveraging the advanced features of AI, we can quickly identify mold, all without the need for professional tools or expertise. Through this project, we hope to demonstrate how AI can transform a common household problem into an easy, and intelligent solution, helping many people live safer and healthier lives}


\subsection{Problem Statement}
\begin{enumerate}
\item Difficulty Identifying Mold 
\item[] Many people struggle to identify mold in its early stages. Especially in apartments and studio housing, residents often notice dark patches on walls, ceilings, or bathrooms but are unsure whether these are mold, dirt, or just discoloration. Without accurate identification, the mold continues to spread, causing damage to the home and potential health issues such as allergies or respiratory irritation. There are currently not a lot of widely known simple and accessible tools for ordinary users to detect mold accurately using only their smartphone. \\
\item Limited Access to Reliable Guidance  
\item[] Even when mold is detected, many people are unsure how to handle it correctly, whether to clean it themselves, what cleaning agents to use, or when to seek professional intervention. Searching online often gives inconsistent or unreliable advice. This causes repeated mold problems and unnecessary expenses for cleaning services or repairs. \\
\item Need for a Simple and Interactive Solution 
\item[] Current mold prevention solutions are often hardware-centric and expensive, such as dehumidifiers or sensors, and a reactive approach, which is cleaning after mold appears. There is a gap in the market for a software-only solution that supports users in all stages: photo-based mold identification, and guidance via chatbot. A tool with AI-driven systems that combine all these features would be more affordable, easier to use, and could help reduce mold-related damage and health issues across many households. 
\end{enumerate}

\subsection{Research on Related Applications}
\begin{enumerate}
\item Mould Identifier: AI Scanner
    \begin{figure}[h!]
    \centering
    \includegraphics[width=0.8\linewidth]{Images/fig1.png}
    \caption{Mould Identifier: AI Scanner}
    \label{fig:enter-label}
\end{figure}

\noindent Mould Identifier is an AI-powered mobile application designed to facilitate easy detection and analysis of mold. By capturing or uploading a photo, the app identifies the mold type, confidence score, assesses associated risks, and provides safety guidance, including a percentage-based indication of severity, danger, exposure, and urgency. The application also features additional informational pages that offer detailed descriptions of different mold types and step-by-step usage guides for effective monitoring and prevention. A history log enables users to track past scans, while the user-friendly interface ensures effortless navigation and accessibility for casual users. \\

\item Mold Identifier: AI Scan
\begin{figure}[h!]
    \centering
    \includegraphics[width=0.8\linewidth]{Images/fig2.png}
    \caption{Mold Identifier: AI Scan}
    \label{fig:enter-label}
\end{figure}

\noindent Mold Identifier is an AI-powered application that detects mold and provides detailed information about the identified species. Upon opening the app, the camera interface launches automatically, allowing users to capture or upload an image for analysis. The AI recognizes the mold type and delivers comprehensive descriptions, including characteristics, typical habitats, and care information. \\

\item Mold Finder AI, Inspection
\begin{figure}[h!]
    \centering
    \includegraphics[width=0.8\linewidth]{Images/fig3.png}
    \caption{Mold Finder AI, Inspection}
    \label{fig:enter-label}
\end{figure}

\noindent Mold Finder AI provides a solution for detecting, analyzing, and understanding mold directly from a mobile device. Utilizing advanced AI technology, the app allows users to capture a photo or upload an existing image, after which it automatically identifies the mold type, analyzes the image, and specifies the location of the mold within the captured area, along with a confidence score. In addition to detection, the app delivers comprehensive analyses covering color, texture, health and safety risks, growth and spread patterns, and recommended remediation measures. \\

\item Virtual Mold Inspection
\begin{figure}[h!]
    \centering
    \includegraphics[width=0.8\linewidth]{Images/fig4.png}
    \caption{Virtual Mold Inspection}
    \label{fig:enter-label}
\end{figure}

\noindent Virtual Mold Inspection is an application platform that enables property owners or homeowners to schedule appointments with certified mold professionals for remote, real-time inspections. Users can select a convenient day and time, then conduct a video call through the app, during which the professional inspector guides the user in examining areas of concern, evaluates potential mold issues, and provides tailored advice on prevention and remediation. Unlike AI-based applications, this platform relies on human expertise rather than automated analysis, which may limit the speed and scalability of assistance. However, the involvement of trained professionals ensures higher accuracy in identifying mold presence and assessing associated risks. \\

\item MoldAI
\begin{figure}[h!]
    \centering
    \includegraphics[width=0.8\linewidth]{Images/fig5.png}
    \caption{MoldAI}
    \label{fig:enter-label}
\end{figure}

\noindent MoldAI is an artificial intelligence–based mold detection application integrated with YCM hardware to provide real-time analysis of mold presence across various surfaces, including walls, shoes, fabrics, and industrial materials. The system incorporates a camera interface that guides users to align the target area within a designated box outline and set the shooting location to ensure consistent image capture. Upon taking or uploading an image, the AI model processes the visual data and classifies whether the area contains mold or not, delivering instant diagnostic feedback. The application also features a dashboard where users can review stored results and access a curated news section that links to external informational resources. However, the current implementation indicates that the detection functionality remains in a prototype phase, with limited real-time analytical performance. \\

\end{enumerate}

\section{Requirement Analysis}
\begin{enumerate}
    \item {\it Common Features}
    \begin{enumerate}[label=(\alph*)]
        \item {\bf Sign-Up and Log-In}
            \begin{itemize}
             \item Enter email: Used for login and account identification. 
             \item Enter password: Must be at least six characters long for security.
             \item Enter username: Used as the display name within the app. Once registered, users can log in using their email and password to access all AI-powered functionalities and save mold detection history. \\
            \end{itemize}
    \end{enumerate}
     
    \item {\it User Specific Features}
        \begin{enumerate}[label=(\alph*)]
        \item {\bf Mold Photo Detection} \\
        Users can upload a photo or take a picture using their device's camera to check whether mold is present using AI image detection. The following are the information that will be displayed: 
            \begin{itemize}
             \item Detection Result 
             \item Location (entered manually by the user for better customization) 
             \item Timestamp
            \end{itemize}
        \item {\bf Home Dashboard (My Home):} \\
        The main screen displays the history, chatbot icon, settings icon and camera icon.
        \item {\bf History Log:} \\
        The application maintains a detection history archive where users can review previous AI analysis results in a chronological list. The history displays the location and the previous uploaded photos with the results. 
        \item {\bf AI Chatbot Assistant:} \\
        The integrated chatbot allows users to ask more information on mold-related questions, such as “What specific products should I use to clean this?”. It leverages a simple NLP model to generate contextual, easy-to-understand advice and explanations.  
        \item {\bf Settings:} \\
        The Settings section includes User Profile and About GomMold. In the User Profile section, users can customize their username and change their password. In the About GomMold section, information about the application is provided.
        \end{enumerate}
\end{enumerate}

\section{Software Development Platform}
\subsection{Development Platform}
    \begin{itemize}
    \item MacOS \\
    MacOS is the chosen and primary chosen development platform as it has software engineering operations like version control management and executing Python script for AI model training are made-easier by its robust UNIX-based interface. Moreover, it is perfect for creating and merging machine learning models due to its excellent support for AI frameworks like TensorFlow and PyTorch. Also, MacOS provides a reliable platform and smooth environment to unify the software engineering and AI components for this project.
    \end{itemize}
    \begin{itemize}
    \item Windows \\
    Windows is a popular and versatile operating system, which makes it a great platform for this project’s development. Windows user-friendly interface simplifies environment configuration and project management, while its robust community support ensures quick and fast. Moreover, this project may use dedicated GPUs for quicker model training because of its strong hardware support. In summary, Windows offers a complete and adaptable environment for creating GomMold application's cross-platform components.
    \end{itemize}

\subsection{Tools and Language}
    \begin{itemize}
        \item Dart \\
        Dart is the programming language used by Flutter, developed and maintained by Google. It is a modern, object-oriented, and type-safe language designed for building fast, reliable, and visually appealing applications across multiple platforms such as Android, iOS, web, and desktop. Dart compiles to native code for mobile and desktop, which ensures smooth performance and quick UI rendering. In this project, Dart is used to develop the frontend mobile application, enabling interactive user experiences and seamless navigation between features such as photo upload, weather display, and AI-based mold detection results. Dart also acts as the bridge connecting the app to backend services, including Firebase for authentication and data storage, and the Flask API for AI model predictions and chatbot communication.\\
        
        \item Firebase \\
        Firebase is a cloud-based Backend-as-a-Service (BaaS) platform developed by Google, offering essential backend tools such as authentication, cloud storage, real-time databases, and analytics. In this project, Firebase plays a central role in managing user accounts, detection history, and image storage. Firebase Authentication secures user registration and login, Cloud Firestore stores detection results and environmental data, and Firebase Storage maintains uploaded mold images. The platform’s scalability and integration with Flutter allow real-time synchronization between the user interface and backend, ensuring that users can access their analysis results and personalized recommendations instantly and securely. \\
        
        \item Flask \\
        Flask is a lightweight and flexible Python-based micro web framework recognized for its simplicity, scalability, and adaptability in modern web development. Its minimalist architecture allows developers to integrate only essential components, promoting high customization, modularity, and efficient system design. Unlike larger frameworks such as Django, Flask offers greater control and faster performance due to its minimal abstraction layers, making it ideal for data-driven and AI-powered applications. In this project, Flask can be used to build RESTful APIs that connect the image analysis and conversational modules into a unified backend system. This implementation enables real-time data exchange, accurate risk assessment, and intelligent feedback, ensuring efficiency and scalability for future smart-home integrations. \\

        \item Flutter \\
        Flutter is an open-source UI software development toolkit created by Google for building visually appealing, cross-platform applications from a single codebase. Using the Dart language, it provides a reactive programming model and a rich collection of pre-designed widgets that allow rapid interface design and customization. Flutter’s hot reload feature supports instant code changes, improving developer productivity and testing efficiency. In this project, Flutter is used to develop the mobile application interface, allowing users to upload photos, view mold detection results, access weather-based risk analysis, and interact with the AI-powered chatbot. Its integration with Firebase and Flask API enables real-time updates and seamless data exchange. \\
        
        \item Javascript \\
        JavaScript is a dynamic and event-driven programming language that enables interactive web functionality and seamless data communication between the client and server. Although it is not used directly in this project’s core development, JavaScript underpins Firebase’s SDKs and cloud functions, which facilitate real-time data synchronization between the mobile app and the backend. Through these SDKs, JavaScript enables secure and efficient operations for user authentication, data storage, and file management without requiring a custom backend. Its asynchronous capabilities allow the system to handle real-time updates such as saving detection results, retrieving history, or synchronizing user preferences, improving the overall responsiveness and scalability of the application. \\

        \item Latex \\
        LaTeX acts as the main tool for preparing this projects' technical documentation. It helps to format paperwork professionally and consistently by following the IEEE report structure accordingly. This project uses an online LaTeX editor like Overleaf to jointly prepare and compile the document within team members. Overleaf’s cloud-based platform enables real-time writing, editing, commenting and document review without version conflicts. Plus, it allows users to do automatic table of contents creation, referencing and also mathematical equations which enhance the user’s quality and readability of the report documentation. So, by using LaTeX in Overleaf, it helps the team to create and polish team project documentations and paperworks. \\
        
        \item OpenAI \\
        OpenAI, founded in 2015, is a leading artificial intelligence research organization dedicated to developing safe and beneficial AI technologies. Its advanced models, such as GPT and DALL·E, provide strong capabilities in natural language understanding and generation, which can be integrated into applications through APIs. In this project, OpenAI is used to power a simple conversational chatbot that assists users by answering mold-related questions and providing general information in a friendly and accessible way. \\
        
        \item Python \\ 
        Python is a high-level, general-purpose programming language known for its simplicity and versatility. It is widely used for backend development, data analysis, machine learning, and automation. In this project, Python serves as the core backend and AI development language, handling the logic for image processing, machine learning, and integration with external APIs. The Flask web framework, written in Python, is used to build the backend API that connects the mobile application to the AI model. Python libraries such as PyTorch and OpenCV are employed for mold detection and image preprocessing, enabling the system to identify potential mold areas and generate predictions based on uploaded images. Additionally, Python facilitates communication with third-party services like OpenAI API for chatbot responses, ensuring that all intelligent features operate smoothly and efficiently within the backend ecosystem. \\
        
        \item Pytorch \\
        PyTorch is one of the primary frameworks used in the development of this project’s AI model. It is a powerful open-source deep learning library that enables efficient training and deployment of neural networks using Python. PyTorch provides an intuitive and flexible platform for building and experimenting with modern AI architectures, including the YOLO model. Its dynamic computation graph allows developers to modify model behavior on the fly, making debugging and experimentation significantly easier. PyTorch’s ecosystem also integrates seamlessly with scientific computing tools and dataset workflows, supporting smooth model development, evaluation, and refinement. In this project, PyTorch plays a central role in training the mold detection model and exporting the final weight file used by the backend for real-time inference. \\

        \item Railway \\
        Railway is a modern cloud deployment platform that allows developers to run applications without managing complex server setups. It provides a simple interface for building, hosting, and scaling applications directly from a GitHub repository, using an automated Linux-based environment. Railway handles the runtime configuration, environment variables, and deployment process, making it much easier and faster compared to traditional cloud hosting. In this project, Railway is used to deploy the Flask backend that powers user authentication, image processing, and machine learning inference. It securely stores API keys, connects seamlessly with Firebase, and provides real-time logs to troubleshoot errors efficiently. By hosting the backend on Railway, the system achieves stable performance, quicker updates, and smooth communication with both the Flutter frontend and the OpenAI chatbot. \\

        \item Roboflow \\
        Roboflow is an essential tool used in this project for dataset creation, management, and preparation. It provides a powerful platform for organizing images, performing annotations, and generating high-quality training datasets for computer vision tasks. Roboflow offers automated preprocessing features such as image resizing, augmentation, and format conversion, which help improve model robustness and performance. Its intuitive web interface makes it easier to handle large and diverse image collections, especially for complex tasks like mold detection where consistent labeling is critical. By exporting datasets directly into formats compatible with YOLO, PyTorch, and ONNX workflows, Roboflow streamlines the entire dataset pipeline. In this project, Roboflow played an important role in curating, augmenting, and exporting the mold dataset, ensuring the model received clean and well-structured data for effective training.\\
        
    \end{itemize}

\subsection{Software in Use}
\begin{itemize}
    \item Figma \\
    \begin{figure}[h!]
    \centering
    \includegraphics[width=0.27\columnwidth]{Images/fig6.png}
    \caption{Figma}
    \label{fig:enter-label}
    \end{figure}
    \item[] Figma is a collaborative online design and prototype tool for user interface (UI), user experience (UX), and digital products. It is one of the most popular tools in the design industry due to its on real-time collaboration. It allows multiple team members to work simultaneously. This is particularly advantageous for this project, as all team members can access and contribute to designs seamlessly. Figma enables users to see the prototype and test the app flow before any code is written. It also provides a dedicated workspace for developers to inspect design elements, extract assets, and translate the design directly into code. \\

    \item Github
    \begin{figure}[h!]
    \centering
    \includegraphics[width=0.35\columnwidth]{Images/fig7.png}
    \caption{Github}
    \label{fig:enter-label}
    \end{figure}
    \item[] Github is the primary code management and version control in this project. It enables team members to work on code independently without overwriting each other’s work. Besides that, it allows to organize branches for different functionality, manage project repositories and as well as making sure all the code changes are recorded with commits. It also prevents data loss in the event of accidental deletion of local files by acting as a secure online backup. Vs code also integrates well with GitHub by offering effortless synchronization between code editing and version control. This project utilizes Github to ensure effective collaboration, accountability and transparency across all phases of the project's development’s. \\

    \item Google Colab
    \begin{figure}[h!]
    \centering
    \includegraphics[width=0.35\columnwidth]{Images/fig12.png}
    \caption{Google Colab}
    \label{fig:enter-label}
    \end{figure}
    \item[] Google Colab serves as the primary model training environment for this project. It provides a cloud-based platform equipped with free GPU and TPU resources, enabling efficient training of the YOLO model without requiring high-performance hardware from team members. Colab offers a collaborative notebook interface that allows code, results, and experiments to be clearly documented and shared. This ensures that training workflows remain reproducible, organized, and easy to update as the dataset evolves. The integration with Google Drive also simplifies data management by enabling seamless loading, saving, and exporting of model weights. Through its accessibility, shared workspace features, and GPU-accelerated environment, Google Colab supports fast experimentation and consistent model development. \\
    
    \item Notion \\
    \begin{figure}[h!]
    \centering
    \includegraphics[width=0.35\columnwidth]{Images/fig8.png}
    \caption{Notion}
    \label{fig:enter-label}
    \end{figure}
    \item[] Notion functions as the primary project management and collaboration platform for the team. It facilitates the organization’s task list,schedules, meeting notes and essential resources within a singular platform. Notion is used to allocate assignments, track the progress and make sure each team member is fully aware of their duties and timelines. Plus, it enables real-time editing and offers simultaneous visibility of updates made by one user to others. By this transparency, it minimizes communication gaps and helps to smoothen the coordination between frontend, backend and AI developer as well. Notion is essential for keeping team alignment and productivity during the development process. \\

    \item Postman
    \begin{figure}[h!]
    \centering
    \includegraphics[width=0.3\columnwidth]{Images/fig11.png}
    \caption{Postman}
    \label{fig:enter-label}
    \end{figure}
    \item[] Postman is an open-source API development and testing tool that allows developers to send requests, analyze responses, and troubleshoot API efficiently. It supports various request types and includes features such as organizing endpoints and automated testing. In this project, Postman is used as the central platform for validating and testing all API interactions before they are integrated into the system. By using Postman as an API testing environment, this project ensures the backend functions correctly.\\
    
    \item Visual Studio Code
    \begin{figure}[h!]
    \centering
    \includegraphics[width=0.35\columnwidth]{Images/fig9.png}
    \caption{Visual Studio Code}
    \label{fig:enter-label}
    \end{figure}
    \item[] Visual Studio Code (VS Code) is a lightweight, open-source code editor developed by Microsoft. It supports multiple programming languages and includes features such as syntax highlighting, debugging, and a wide range of plugin extensions. This project uses VS Code as the main platform for all stages of code development. It serves as the central hub where both front-end and back-end code is written, tested, and managed before being deployed to the physical application. By using VS Code as the primary development environment, this project can seamlessly integrate all components, ensuring that the final project is well-organized and easily maintainable.
    
\end{itemize}

\subsection{Team's Development Environment}

\begin{table}[htp]
 \centering
 \setlength{\tabcolsep}{12pt}
 \renewcommand{\arraystretch}{1.3}
 \begin{tabular}{| >{\centering\arraybackslash}m{2.3cm} | >{\centering\arraybackslash}m{4cm} |}
 \hline
 {Name} & {Environment} \\ 
 \hline

 Aina Nur Insyeerah & 
 \begin{tabular}[c]{@{}c@{}}
 Macbook Air M4, 24GB \\ 
 MacOS Tahoe 26.0.1 \\ 
 Python 3.13.7
 \end{tabular} \\ 
 \hline

 Nur Ain Syafiqah &
 \begin{tabular}[c]{@{}c@{}}
 Macbook Air M3, 24GB \\ 
 MacOS Sequoia 15.6.1 \\ 
 Python 3.9.6
 \end{tabular} \\
 \hline

 Nur Shaqeerah &
 \begin{tabular}[c]{@{}c@{}}
 Macbook Air M1, 8GB \\ 
 MacOS Sequoia 15.6.1 \\ 
 Python 3.13.7
 \end{tabular} \\
 \hline

 Yumni Karmila &
 \begin{tabular}[c]{@{}c@{}}
 HP Pavillion 15, \\
 12th Gen Intel(R) Core(TM) \\
 i7-1255U (1.70 GHz), 512GB, \\
 16 RAM, 2GB GPU, Python 3.10.18
 \end{tabular} \\
 \hline
 
 \end{tabular}
\end{table}

\subsection{Cost Estimation}
This project incurs no additional financial cost, as it relies entirely on open-source and free software. Development is carried out using personal laptops, eliminating hardware expenses. Additionally, the project takes advantage of cloud services, such as computing and storage resources, which provide scalable and flexible infrastructure without extra cost.

\subsection{Task Distribution}

\begin{table}[h]
\centering
\setlength{\tabcolsep}{12pt}
\renewcommand{\arraystretch}{1.5}
\begin{tabular}{|p{1cm}|p{1.5cm}|p{4cm}|}
\hline Tasks & Name & Descriptions\\ \hline
Frontend Developer & Nur Ain Syafiqah, Yumni Karmila & Frontend developers build the user interface, manage interactivity, and integrate real-time data using Flutter and Dart. The first developer, the UI implementation Specialist transforms Figma designs into responsive mobile interfaces, handling navigation, animations, and Firebase authentication with focus on usability and performance. The second developer is the Data Integration and Visualization Specialist who connects the Flask backend with the frontend through APIs for mold detection, manages data visualizations, synchronization and collaborates with backend and AI teams for seamless functionality.\\
\hline
UI-UX Designer & Yumni Karmila & The UX Designer shapes the overall user experience and visual design of the app, ensuring accessibility, usability,and consistency. Using Figma, the designer create interface layouts, color schemes, typography, and user flows through wireframes and interactive prototypes. Usability testing is conducted to refine designs based on feedback and collaborate with frontend developers to ensure accurate visual implementation. The \\
\hline
\end{tabular}
\end{table}

\begin{table}[h]
\centering
\setlength{\tabcolsep}{12pt}
\renewcommand{\arraystretch}{1.5}
\begin{tabular}{|p{1cm}|p{1.5cm}|p{4cm}|}
\hline
&& designer also maintains a design system for consistent branding.\\
\hline
AI Developer & Nur Shaqeerah & The AI Developer is responsible for deploying a YOLO-based machine learning model to detect and classify mold in real time. This includes collecting and preprocessing image data, training and fine-tuning the YOLO model in Google Colab using PyTorch for higher accuracy, and integrating the trained model into the app through Flask API endpoints for seamless communication between the AI, backend, and frontend. By leveraging YOLO’s real-time detection capabilities, this role ensures automated, scalable, and reliable mold recognition.\\
\hline

Backend Developer & Aina Nur \newline Insyeerah & The Backend Developer designs the database, builds APIs, and ensures smooth data flow between the frontend, backend, and AI components. Using Python and Flask, the developer creates a robust server system that processes user data and returns real-time responses. The role integrates external APIs, uses an ONNX-based cloud model for mold detection, implements the OpenAI API for chatbot features, and manages secure data storage with Firebase. The backend is deployed on Railway to ensure stable, scalable, and efficient cloud performance. \\
\hline

\end{tabular}
\end{table}

 \newpage

\section{Specifications}
\subsection{Initial Page (Splash Screen)}
\begin{figure}[h!]
    \centering
    \includegraphics[width=0.35\columnwidth]{Images/Initial page.png}
    \caption{Initial page}
    \label{fig:enter-label}
    \end{figure}
\noindent The initial screen appears for a few seconds upon launching the application, serving as a splash page that introduces the GomMold brand and ensures all required system resources are initialized (e.g., Firebase, API connections).
\begin{itemize}
    \item {\it Process}
    \begin{itemize}
        \item The app loads core dependencies (Flask API endpoints, AI model, Firebase configurations).
        \item Displays GomMold logo and animation for 3–5 seconds.
        \item Automatically transitions to the Log-In Page once initialization completes.
    \end{itemize}
\item {\it System Response}
    \begin{itemize}
        \item Success: Redirects to Log-In Page.
        \item Failure: Displays “Connection Error” message if backend initialization fails.
    \end{itemize}
\end{itemize}

\subsection{Log In}
\begin{figure}[h!]
    \centering
    \includegraphics[width=0.35\columnwidth]{Images/LogIn Page.png}
    \caption{Log In page}
    \label{fig:enter-label}
    \end{figure}
\noindent This page allows registered users to access their account using a valid email and password.
    \begin{itemize}
        \item Input Fields:
        \begin{itemize}
            \item Email: Used for account identification (case-insensitive).
            \item Password: Must match with the securely hashed password stored in the database.
        \end{itemize}
        \item Buttons:
        \begin{itemize}
            \item Log In: Validates Credentials
            \item Sign Up: Redirects to Sign Up page for new users.
        \end{itemize}
        \item Process Flow:
        \begin{enumerate}
            \item User enters login credentials
            \item Upon pressing “Log In,” the frontend sends a \texttt {POST /api/auth/login} request to the Flask backend.
            \item The backend checks if the email exists in the users collection in Firestore.
            \item If found, it validates the password using check\_password\_hash.
            \item On successful authentication, a JWT token (valid for 24 hours) is generated and returned to the app.
            \item User is redirected to the Homepage.
        \end{enumerate}
\item {\it System Response}
    \begin{itemize}
        \item Success: Redirects to Homepage.
        \item Failure: Displays "Login failed. Try again."\\
    \end{itemize}
\end{itemize}

\subsection{Sign Up}
\begin{figure}[h!]
    \centering
    \includegraphics[width=0.35\columnwidth]{Images/Sign Up Page.png}
    \caption{Sign Up page}
    \label{fig:enter-label}
    \end{figure}
\noindent This page allows new users to register an account.
    \begin{itemize}
        \item Input Fields:
        \begin{itemize}
            \item Email: Must be unique in the database
            \item Username: Must be unique in the database.
            \item Password: Minimum of 6 characters.
        \end{itemize}
        \item Buttons:
        \begin{itemize}
            \item Sign Up: Registers new user.
            \item Back to Log In: Returns to Log-In Page.
        \end{itemize}
        \item Process Flow:
        \begin{enumerate}
            \item User enters email, username, and password.
            \item Frontend sends a \texttt {POST /api/auth/signup} request to Flask backend.
            \item Backend checks if email or username already exist in the users collection.
            \item If unique, the password is hashed using Werkzeug’s generate\_password\_hash() and stored securely.
            \item Upon successful registration, a JWT token is generated and returned.
            \item The user is redirected to the Log-In Page.
        \end{enumerate}
\item {\it System Response}
    \begin{itemize}
        \item Success: Displays “Registration successful. Please log in."
        \item Failure: Displays one of the following:
        \begin{itemize}
            \item "Password must be at least 6 characters."
            \item "Sign up failed. Try again."
        \end{itemize}
    \end{itemize}
\end{itemize}

\subsection{Homepage (Main Dashboard)}
\begin{figure}[h!]
    \centering
    \includegraphics[width=0.35\columnwidth]{Images/Homepage.png}
    \caption{Homepage}
    \label{fig:enter-label}
    \end{figure}
\noindent The Homepage serves as the user’s control center. It displays detection history and provides quick navigation to camera, chatbot, and settings.
    \begin{itemize}
        \item Display Elements:
        \begin{itemize}
            \item Middle Section (History): Displays list of past detections:
            \begin{itemize}
                \item Location
                \item Results (Mold / No Mold)
            \end{itemize}
        \end{itemize}
        \item Buttons Navigation Bar:
        \begin{itemize}
            \item Settings Button (Left) – Opens Settings Page.
            \item Camera Button (Center) – Opens Image Page.
            \item Chatbot Button (Right) – Opens Chatbot Page.
        \end{itemize}
        \item Process Flow:
        \begin{enumerate}
            \item When the page loads, it sends \texttt {GET /api/mold/history} to the backend.
            \item Backend retrieves the user’s history from the detections table and sends it to the frontend.
            \item Backend checks if username exists in the users table.
            \item Data is displayed in chronological order.
        \end{enumerate}
\item {\it System Response}
    \begin{itemize}
        \item Success: Displays recent detection history.
        \item Displays “No detection history found.”\\
    \end{itemize}
\end{itemize}

\subsection{Image Page}
\begin{figure}[h!]
    \centering
    \includegraphics[width=0.35\columnwidth]{Images/Identify page.png}
    \caption{Image page}
    \label{fig:enter-label}
    \end{figure}
\noindent This page allows users to choose between capturing a photo or uploading one from their gallery for AI-based mold detection.
    \begin{itemize}
        \item Buttons:
        \begin{itemize}
            \item Take Photo: Opens real-time camera interface.
            \item Choose from Gallery: Opens device’s photo picker.
            \item Back: Returns to Homepage.
        \end{itemize}
        \item Process Flow:
        \begin{enumerate}
            \item When “Take Photo” is selected, app requests camera permission.
            \item When “Choose from Gallery” is selected, app opens local gallery.
            \item Once an image is chosen or captured, it navigates to the Identify` Page.
        \end{enumerate}
\item {\it System Response}
    \begin{itemize}
        \item Success: Redirects to Identify Page with selected image.
        \item Failure: Displays “Access to camera denied” if camera access is blocked.\\
    \end{itemize}
\end{itemize}

\subsection{Identify Page (Mold Detection Page)}

\begin{figure}[h!]
    \centering
    \includegraphics[width=0.35\columnwidth]{Images/Identify Page - Success Result (1).png}
    \includegraphics[width=0.35\columnwidth]{Images/Identify Page - Failed Result (1).png}
    \caption{Mold Detection Page}
    \label{fig:enter-label}
\end{figure}

\noindent The Identify Page allows users to capture or upload an image and run mold detection using the backend AI model. 
Results are displayed immediately with a color-coded status to indicate safety.

\begin{itemize}
    \item Buttons:
    \begin{itemize}
        \item Live Camera (real-time capture)
        \item Shoot Button
        \item Back Button (returns to Image Page)
        \item Result Display Area
    \end{itemize}

    \item Process Flow:
    \begin{enumerate}
        \item User takes or selects an image from the device.
        \item The image is sent to the backend via 
              \texttt{POST /api/mold/detect} along with the user's authentication token.
        \item Flask receives the uploaded file, stores it in Firebase Storage, and generates a public image URL.
        \item The backend temporarily saves the file and passes it into the ONNX model using \texttt{predict\_image()}.
        \item The model outputs bounding boxes and confidence scores for potential mold.
        \item Based on the filtered detections, the backend determines:
        \begin{itemize}
            \item Red Warning – "Mold Detected" (if at least one prediction is above threshold)
            \item Green Safe – "No Mold Detected" (if no valid detections are found)
        \end{itemize}
        \item The backend saves the result, predictions, image URL, and timestamp into Firestore under the \texttt{detections} collection.
        \item The JSON response is returned to the app, containing the result, prediction details, and timestamp.
        \item The app displays the detection result instantly on the same page.
    \end{enumerate}

    \item \textbf{System Response}
    \begin{itemize}
        \item Success: Detection result is displayed with prediction data.
        \item Failure: Shows ``Image processing error.''
    \end{itemize}
\end{itemize}


\subsection{Chatbot Page}
\begin{figure}[h!]
    \centering
    \includegraphics[width=0.35\columnwidth]{Images/Chatbot.jpg}
    \caption{Chatbot page}
    \label{fig:enter-label}
    \end{figure}
\noindent The chatbot allows users to interact with an AI assistant for mold-related questions and general information.

\begin{itemize}

    \item Components:
    \begin{enumerate}
        \item Initial AI message introducing itself.
        \item Text input field and Send button.
        \item “Back” button (returns to Homepage).
    \end{enumerate}

    \item Process Flow
    \begin{enumerate}
        \item Upon entering, the chatbot sends a greeting via backend API: \texttt{GET /api/chatbot/start}
        \item User types a question into the chat input.
        \item The query is sent to the backend: \texttt{POST /api/chatbot/query}
        \item Backend forwards the message to the OpenAI API for language processing.
        \item The AI generates a response and returns it to the app.
    \end{enumerate}

    \item \textit{System Response}
    \begin{itemize}
        \item Success: AI reply is displayed in a chat-style conversation.
        \item Failure: Displays “Chat service unavailable.”
    \end{itemize}

\end{itemize}

\item {\it System Response}
    \begin{itemize}
        \item Success: AI reply displayed in conversation format.
        \item Failure: Displays “Chat service unavailable.”
    \end{itemize}
\end{itemize}

\subsection{Settings Page}
\begin{figure}[h!]
    \centering
    \includegraphics[width=0.35\columnwidth]{Images/Settings.png}
    \caption{Settings page}
    \label{fig:enter-label}
    \end{figure}
\noindent This page allows users to manage their profile information and view app details.
    \begin{itemize}
        \item Display Elements:
        \begin{itemize}
            \item Change Username Button
            \item Change Password Button
            \item About GomMold Section (App details, version info, developer credits)
            \item Back Button: Returns to Homepage
        \end{itemize}
        \item Process Flow
        \begin{enumerate}
            \item When user clicks “Change Username” or “Change Password,” a form opens.
            \item On submission, app sends a \texttt{PATCH /api/users/update} request.
            \item Backend updates user data in the users table.
            \item Updated info is synced with Firebase Authentication if applicable.
        \end{enumerate}
\item {\it System Response}
    \begin{itemize}
        \item Success: Displays “Profile updated successfully.”
        \item Displays “Error updating profile.”
    \end{itemize}
\end{itemize}

\section{Architecture page}
\subsection{Overall Architecture}
\begin{figure}[h!]
    \centering
    \includegraphics[width=1.0\linewidth]{Images/architecture diagram.png}
        \caption{Application Architecture}
    \label{fig:enter-label}
\end{figure}

The GomMold system architecture is organized into two primary modules: Frontend and Backend. The Frontend module, developed using Flutter, serves as the main interface through which users interact with the application. It provides functionalities such as image upload, mold analysis viewing, detection history, and chatbot communication. The user interface emphasizes accessibility and responsiveness, allowing users to capture or select images, review results, and manage their profile settings. The frontend collects all user inputs and communicates with the backend through secure RESTful APIs, relying entirely on the backend to perform authentication, verification, and session handling.\\

The Backend module, implemented with the Flask framework and deployed on a Linux-based server via Railway, acts as the central processing engine of the system. It manages all authentication processes, including user registration, login, token generation, and session validation, ensuring secure access across devices. The backend also handles the core image-processing pipeline: receiving uploaded images, running them through the machine learning inference engine, interpreting prediction outputs, and returning structured JSON responses back to the Flutter application. Additionally, the backend integrates the OpenAI API to power an interactive chatbot, enabling users to ask mold-related questions and receive cleaning advice, risk explanations, and preventative guidance. By hosting the backend on Railway, GomMold benefits from faster deployment cycles, simpler configuration, and more transparent error handling.\\

The data storage layer is responsible for managing all user-generated and application-related data. Firebase Firestore is used to store structured data such as user profiles and detection metadata, while Firebase Storage handles unstructured data such as uploaded images and model inference results. This design ensures efficient data retrieval, real-time synchronization, and reliable image versioning across the application layers.\\

The machine learning component forms the core analytical engine of GomMold. The mold detection model is built using PyTorch and trained offline using datasets prepared and annotated through Roboflow. Training is conducted in Google Colab to optimize the model under various real-world lighting and surface conditions. After training, the model is exported to ONNX format to ensure compatibility with the cloud backend. Within the Flask backend, the ONNX model is loaded and executed using ONNX Runtime, allowing fast and lightweight inference for every uploaded image. The system evaluates the model’s output to determine the presence of mold and its confidence level, returning the result directly to the user through the mobile application. In addition to detection, GomMold incorporates an optional chatbot powered by the OpenAI API, enabling users to obtain further clarification, cleaning tips, or preventive advice.\\

Collectively, these modules create an integrated ecosystem that combines real-time image analysis, cloud-based data management, and AI-assisted feedback to deliver a practical, reliable, and user-friendly solution for mold detection and prevention. The modular architecture ensures scalability, maintainability, and interoperability, positioning GomMold as a viable tool for both everyday users and future advancements in environmental health monitoring.


\subsection{Directory Organization} 
\begin{table}[h]
\caption{DIRECTORY-ORGANIZATION-FRONTEND-1}
\def\arraystretch{1.24} \small
    \begin{tabular}{|p{3.7cm}|p{4.1cm}|}
\hline
Directory & File Name \\ \hline
          Root  & .flutter-plugins-dependencies \par .gitignore \par .metadata \par README.md \par analysis\_options.yaml \par devtools\_options.yaml \par firebase.json \par gomMold\_frontend.iml \par package-lock.json \par package.json \par pubspec.lock \par pubspec.yaml
          \\ \hline
\end{tabular}
\end{table}
\newpage

          
\begin{table}[h]
\caption{DIRECTORY-ORGANIZATION-FRONTEND-2}
\def\arraystretch{1.24} \small
    \begin{tabular}{|p{3.7cm}|p{4.1cm}|}
\hline
          
          .dart\_tool \par   & \par package\_config.json \par package\_graph.json \par version 
          \\ \hline
          
          .dart\_tool/dartpad \par  &\par web\_plugin\_registrant.dart
          \\ \hline
          
          .idea \par   & modules.xml \par workspace.xml
          \\ \hline
          
          .idea/libraries \par & Dart\_SDK.xml \par KotlinJavaRuntime.xml
          \\ \hline
          
          .idea/runConffigurations \par  & main\_dart.xml  
          \\ \hline

          android \par & .gitignore \par build.gradle.kts \par gommold\_frontend\_android.iml \par gradle.properties \par gradlew \par gradlew.bat \par local.properties \par settings.gradle.kts
          \\ \hline
          
          android/app \par   & \par build.gradle.kts \par google-services.json \par 
          \\ \hline

          android/app/src/debug \par   & \par AndroidManifest.xml
          \\ \hline

          android/app/src/main \par   & \par AndroidManifest.xml
          \\ \hline

          android/app/src/debug \par   & \par AndroidManifest.xml
          \\ \hline

          android/app/src/main/kotlin/ \par com/example/\par gommold\_frontend  \par & \par MainActivity.kt
          \\ \hline

          android/app/src/main/res \par /drawable-v21 \par   & \par launch\_background.xml
          \\ \hline

          android/app/src/main/res \par /drawable \par   & \par launch\_background.xml
          \\ \hline

          android/app/src/main/res  \par /mipmap-hdpi \par   & \par ic\_launcher.png
          \\ \hline

          android/app/src/main/res \par /mipmap-mdpi \par   & \par ic\_launcher.png
          \\ \hline

          android/app/src/main/res \par /mipmap-xhdpi \par   & \par ic\_launcher.png
          \\ \hline

          android/app/src/main/res  \par /mipmap-xxhdpi \par   & \par ic\_launcher.png
          \\ \hline
          
          android/app/src/main/res  \par /mipmap-xxxhdpi \par   & \par ic\_launcher.png
          \\ \hline

          
\end{tabular}
\end{table}

\newpage 
\begin{table}[h]
\caption{DIRECTORY-ORGANIZATION-FRONTEND-3}
\def\arraystretch{1.24} \small
    \begin{tabular}{|p{3.7cm}|p{4.1cm}|}
\hline
          
          android/app/src/main/res  \par /values-night \par   & \par styles.xml
          \\ \hline

          android/app/src/main/res  \par /values \par   & \par styles.xml
          \\ \hline
          
          android/app/src/profile \par   & \par AndroidManifest.xml
          \\ \hline

          android/gradle/wrapper \par   & \par gradle-wrapper.properties
          \\ \hline

          assets/images \par   & \par ChatBot.png \par Homepage-time.png \par Identify Page- \par Success Result.png  \par Image Page.png \par LogIn Page.png \par Settings.png \par Sign Up Page.png \par logo.png  
          \\ \hline
          
          dist   & \par socket.js
          \\ \hline

          ios \par   & \par .gitignore \par Podfile
          \\ \hline

          ios/Flutter \par   & \par AppFrameworkInfo.plist \par Debug.xconfig \par Generated.xconfig  \par Release.xconfig \par flutter\_export\_environment.sh
          \\ \hline

          ios/Runner.xcodeproj \par   & \par project.pbxproj
          \\ \hline

          ios/Runner.xcodeproj \par /project.xcworkspace   & \par contents.xcworkspacedata
          \\ \hline

          ios/Runner.xcodeproj \par /project.xcworkspace \par /xcshareddata  & \par IDEWorkspaceChecks.list
          \\ \hline

          ios/Runner.xcodeproj \par /xcshareddata/xcschemes  & \par Runner.xscheme
          \\ \hline

          ios/Runner.xcworkspace \par   & \par contents.xcworkspacedata
          \\ \hline

          ios/Runner.xcworkspace \par /xcshareddata  & \par IDEWorkspaceChecks.list \par WorkspaceSettings.xcsettings
          \\ \hline

          ios/Runner \par   & \par AppDelegate.swift \par GeneratedPluginRegistrant.h \par GeneratedPluginRegistrant.m \par GoogleService-Info.plist \par Info.plist \par Runner-Bridging-Header.h
          \\ \hline

           ios/Runner/Assets.xcassets \par /AppIcon.appiconset  & \par Contents.json \par Icon-App-1024x1024@1x.png \par Icon-App-20x20@1x.png \par Icon-App-20x20@2x.png \par Icon-App-20x20@3x.png \par Icon-App-29x29@1x.png \par Icon-App-29x29@2x.png \par Icon-App-29x29@3x.png \par Icon-App-40x40@1x.png \par Icon-App-40x40@2x.png \par Icon-App-40x40@3x.png \par Icon-App-60x60@2x.png \par Icon-App-60x60@3x.png \par Icon-App-76x76@1x.png \par Icon-App-76x76@2x.png \par Icon-App-83.5x83.5@2x.png 
          \\ \hline
          

\end{tabular}
\end{table}

\begin{table}[h]
\caption{DIRECTORY-ORGANIZATION-FRONTEND-4}
\def\arraystretch{1.24} \small
    \begin{tabular}{|p{3.7cm}|p{4.1cm}|}
\hline
        

          ios/Runner/Assets.xcassets \par /LaunchImage   & \par Contents.json
          \\ \hline

          ios/Runner/Assets.xcassets \par /LaunchImage  & \par Contents.json
          \\ \hline

          imageset \par   & \par LaunchImage.png \par LaunchImage@2x.png \par LaunchImage@3x.png  \par REDME.md
          \\ \hline

          ios/Runner/Base.lproj    & \par LaunchScreen.storyboard \par Main.storyboard
          \\ \hline

          ios/RunnerTests  & \par RunnerTests.swift
          \\ \hline

          lib   & \par main.dart
          \\ \hline

          lib/src/models   & \par detection.dart
          \\ \hline

          lib/src/screens   & \par chatboot\_page.dart \par homepage\_time.dart \par identify\_page.dart \par image\_page.dart \par initial\_page.dart \par login\_page.dart \par settings\_page.dart \par sign\_up\_page.dart
          \\ \hline      

          lib/src/services \par   & \par api\_service.dart, auth\_service.dart, firebase\_service.dart 
        \\ \hline
        
        lib/src/utils  & \par constants.dart \par image\_utils.dart
        \\ \hline   

        lib/src/widgets  & \par card\_item.dart \par custom\_button.dart  \par custom\_textfield.dart
        \\ \hline 

        linux  & \par .gitignore \par CMakeLists.txt
        \\ \hline 

        linux/flutter   & \par CMakeLists.txt \par generated\_plugin\_registrant.cc \par generated\_plugin\_registrant.h \par generated\_plugins.cmake
        \\ \hline

        linux/runner \par   & \par CMakeLists.txt \par main.cc \par my\_application.cc \par my\_application.h
        \\ \hline

        macos \par   & \par .gitignore \par Podfile \par Podfile.lock
        \\ \hline

        macos/Flutter \par   & \par Flutter-Debug.xcconfig \par Flutter-Release.xcconfig \par GeneratedPluginRegistrant.swift
        \\ \hline

        macos/Runner.xcodeproj  & \par project.pbxproj
         \\ \hline

        macos/Runner.xcodeproj \par /project.xcworkspace \par /xcshareddata   & \par IDEWorkspaceChecks.plist
        \\ \hline

        macos/Runner.xcodeproj \par /xcshareddata/xcschemes   & \par Runner.xcscheme
        \\ \hline

        macos/Runner.xcworkspace  & \par contents.xcworkspacedata
        \\ \hline

        macos/Runner.xcworkspace \par /xcshareddata  & \par IDEWorkspaceChecks.plist
        \\ \hline

        macos/Runner \par   & \par AppDelegate.swift \par DebugProfile.entitlements \par GoogleService-Info.plist \par Info.plist \par MainFlutterWindow.swift \par Release.entitlements
         \\ \hline
        
    \end{tabular}
\end{table}

\begin{table}[h]
\caption{DIRECTORY-ORGANIZATION-FRONTEND-5}
\def\arraystretch{1.24} \small
    \begin{tabular}{|p{3.7cm}|p{4.1cm}|}
\hline


        macos/Runner/Assets.xcassets \par /AppIcon.appiconset \par   & \par Contents.json \par app\_icon\_1024.png \par app\_icon\_128.png \par app\_icon\_16.png \par app\_icon\_256.png \par app\_icon\_32.png \par app\_icon\_512.png \par app\_icon\_64.png
        \\ \hline

        macos/Runner/Base.lproj   & \par MainMenu.xib
        \\ \hline

        macos/Runner/Configs \par   & \par AppInfo.xcconfig \par Debug.xcconfig \par Release.xcconfig \par Warnings.xcconfig
        \\ \hline

        macos/RunnerTests  & \par RunnerTests.swift
        \\ \hline

        node\_modules   & \par .package-lock.json
        \\ \hline

        node\_modules/ws   & \par LICENSE \par README.md \par browser.js \par index.js \par package.json \par wrapper.mjs
        \\ \hline

        node\_modules/ws/lib \par   & \par buffer-util.js \par constants.js \par event-target.js \par extension.js \par limiter.js \par permessage-deflate.js \par receiver.js \par sender.js \par stream.js \par subprotocol.js \par validation.js \par websocket-server.js \par websocket.js
        \\ \hline

        test   & \par widget\_test.dart
        \\ \hline

        web & \par favicon.png \par index.html \par manifest.json
        \\ \hline

        web/icons \par   & \par Icon-192.png \par Icon-512.png \par Icon-maskable-192.png \par Icon-maskable-512.png
        \\ \hline

        windows \par   & \par .gitignore \par CMakeLists.txt
        \\ \hline

        windows/flutter \par   & \par CMakeLists.txt \par Generated\_plugin\_registrant.cc \par Generated\_plugin\_registrant.h \par generated\_plugin.cmake
        \\ \hline

        windows/runner \par   & \par CMakeLists.txt \par Runner.rc \par flutter\_window.cpp \par flutter\_window.h \par main.cpp \par resource.h \par resource.h \par runner.exe.manifest \par utils.cpp \par utils.h \par wiin32\_window.cpp \par win32\_window.h
        \\ \hline

        windows/runner/resources   & \par app\_icon.ico
        \\ \hline
        
\end{tabular}
\end{table}

\begin{table}[h!]
\caption{DIRECTORY-ORGANIZATION-BACKEND-1}
\def\arraystretch{1.24} \small
    \begin{tabular}{|p{3.8cm}|p{4.3cm}|}
        \hline
        Directory & File Name \\ \hline
          root \par   & \par .gitignore \par Procfile \par README.md \par requirements.txt \par runtime.txt
        \\ \hline

        src \par   & \par \_\_init\_\_.py \par app.py \par firebase\_init.py \par ml\_model.py \par token\_utils.py
        \\ \hline

        src/routes \par   & \par \_\_init\_\_.py \par auth.py \par chatbot.py \par history.py \par mold.py \par user.py
        \\ \hline

        src/models \par   & \par best.onnx
        \\ \hline

\end{tabular}
\end{table}

\begin{table}[h]
\caption{DIRECTORY-ORGANIZATION-AI}
\def\arraystretch{1.24} \small
    \begin{tabular}{|p{3.9cm}|p{4cm}|}
    \hline
        Directory & File Name \\ \hline
        root \par   & \par .DS\_Store \par README.md
        \\ \hline

        models \par   & \par README.md (download link for the trained model because github does not allow large file upload)
        \\ \hline

        training \par   & \par GomMold\_ML.ipynb \par dataset\_info.md
        \\ \hline

        training/images \par   & \par BoxF1\_curve (1).png \par BoxPR\_curve (1).png \par README.md \par confusion\_matrix (1).png
        \\ \hline
        
    \end{tabular}
\end{table}

\clearpage
\subsection{Module 1: Frontend}
\begin{enumerate}
    \item[1)] Purpose\\
    The GomMold Frontend Module is the main user-facing interface, managing all visual elements, inputs, and navigation to ensure an intuitive, responsive experience across devices. Built with Flutter, it securely communicates with the backend to upload images, receive mold detection results, and handle user authentication. Its goal is to provide a smooth, accessible, and reliable platform for efficient mold identification while maintaining consistent performance and usability.\\ 
    \item[2)] Functionality \\
    The GomMold frontend provides full user-facing functionality, including authentication, secure token storage, and input handling. Users can take photos or choose images from the gallery, and these are sent to the backend through RESTful APIs for analysis. Detection results are displayed clearly, and navigation between pages such as login, home, and identification is smooth. The interface uses reusable widgets for consistent design and manages state with provider and service classes for real-time interaction with the backend. The frontend also includes a chatbot, data visualization, history viewing, and profile management, offering a complete and user-friendly mold detection experience. \\
    
    \item[3)]Location of source code : \url{https://github.com/GomMold/GomMold_FrontEnd}\\ 
    
\end{enumerate}

\subsection{Module 2 : Backend}
\begin{enumerate}
    \item[1)] Purpose\\
  The GomMold backend serves as the system’s central processing and data-management layer. Built with Python and Flask, it handles user requests, authentication, and secure data exchange between the mobile app, AI model, and cloud storage. It manages user data, stores images, and enables real-time communication with the detection model. Deployed on Railway, the backend delivers a stable and scalable environment for the system’s core functions.\\
    \item[2)] Functionality\\
   The GomMold backend handles all core server functions, including user authentication, image processing, and communication with the AI model for mold detection. It returns predictions with mold types and confidence scores, and integrates the OpenAI API to provide chatbot-based cleaning and prevention advice. User data, images, and detection history are stored using Firebase Firestore and Google Cloud Storage. All modules communicate through encrypted API endpoints, ensuring secure and reliable operation. Overall, the backend delivers the infrastructure supporting GomMold’s cloud-based intelligent mold detection system.\\
    \item[3)]Location of source code : \url{https://github.com/GomMold/GomMold_BackEnd}

\end{enumerate}
\subsection{Module 3 : AI}
\begin{enumerate}
\item[1)] Purpose\\
  The GomMold AI Module is the system’s core analytical engine, performing automated mold detection and classification on user images. Using a YOLO-based deep learning model in PyTorch, it delivers accurate, real-time recognition. The module covers data collection, preprocessing, training, and fine-tuning to ensure strong real-world performance. Once trained, the model is deployed through the Flask backend for seamless system integration. It also supports optional OpenAI-powered conversational features that provide cleaning and prevention guidance. Its primary goal is to offer reliable, scalable, and efficient mold detection that aids user decision-making and enhances the overall application.\\
\item[2)] Functionality\\
   The GomMold AI Module receives images uploaded through the frontend, preprocesses them, and executes inference using the trained YOLO model to determine mold presence and associated confidence scores. Image datasets are annotated and prepared via Roboflow, trained in Google Colab, and deployed for real-time analysis within the backend. Detection results are transmitted back to the frontend through secure API endpoints, ensuring clear and interpretable outcomes for users. Furthermore, the module supports an optional OpenAI-powered chatbot that provides general cleaning tips, preventive guidance, and interactive follow-up responses. Overall, the AI Module ensures efficient AI-driven analysis, seamless integration with other system components, and an enhanced, data-informed user experience.\\
   
\item[3)] Custom Mold Detection Dataset (Roboflow – v4)\\
\noindent The dataset contains 1,134 annotated images of household surfaces.
    \begin{itemize}
        \item It consists of two classes:
            \begin{itemize}
                \item Mold
                \item No Mold
            \end{itemize}
            
        \item Images were collected from:
            \begin{itemize}
                \item Personal photography
                \item Additional samples from public Roboflow Universe datasets
                \item Additional samples from google 
            \end{itemize} 
            
         \item Roboflow was used for:
            \begin{itemize}
                \item Annotation
                \item Dataset augmentation
                \item Resizing all images to 640×640
                \item Exporting to YOLOv8 format\\
            \end{itemize} 
    \end{itemize}
    
\item[4)] System Configuration\\ 
\noindent The training was conducted using Google Colab GPU environment.
    \begin{itemize}
        \item Hardware (Colab) :
            \begin{itemize}
                \item GPU: Tesla T4 
                \item RAM: ~12–15GB
                \item GPU VRAM: 16GB
            \end{itemize}
            
        \item Software :
            \begin{itemize}
                \item Python 3.10
                \item PyTorch (CUDA-enabled)
                \item Ultralytics YOLOv8
                \item Roboflow SDK\\
            \end{itemize} 
    \end{itemize}

\item[5)]Location of source code : \url{https://github.com/GomMold/GomMold_ML}\\
 
\end{enumerate}

\section{Use Case}
\subsection{Initial Page}
\begin{figure}[h!]
    \centering
    \includegraphics[width=0.35\columnwidth]{Images/Rinitialpage.png}
    \caption{Initial page}
    \label{fig:enter-label}
    \end{figure}
\noindent When users open the GomMold app, they are first greeted with a minimal welcome screen featuring the logo and a “Get Started” button. This page serves as the entry point of the application, guiding users smoothly into the main experience. From here, they can proceed to authentication or move directly into the app based on their login status.

\subsection{Sign Up Page}
\newpage
\begin{figure}[h!]
    \centering
    \begin{minipage}{0.35\columnwidth}
        \centering
        \includegraphics[width=\linewidth]{Images/Rsignuperror.png}
        \caption{Sign up error}
        \label{fig:sign up page}
    \end{minipage}%
    \begin{minipage}{0.35\columnwidth}
        \centering
        \includegraphics[width=\linewidth]{Images/Rsignuperror2.png}
        \caption{Sign up error}
        \label{fig:signup page}
    \end{minipage}
\end{figure}
\noindent When users enter incorrect or incomplete information on the Sign Up page, the app provides instant validation feedback by highlighting specific errors such as an invalid email format or a password that is too short. This allows users to correct their input before continuing. If the registration still fails due to invalid details or an email that already exists, the page displays an error message prompting the user to try again, ensuring a clear and guided sign-up process.

\subsection{Login Page}
\begin{figure}[h!]
    \centering
    \includegraphics[width=0.35\columnwidth]{Images/Rloginpage.png}
    \caption{Login Page}
    \label{fig:enter-label}
    \end{figure}
\noindent Users enter their email and password, and the system sends these credentials to the backend for verification. The backend validates the information against existing records, and if the credentials are correct, the user is authenticated and granted access to the application. 
\newpage
\begin{figure}[h!]
    \centering
    \begin{minipage}{0.35\columnwidth}
        \centering
        \includegraphics[width=\linewidth]{Images/Rloginerror.png}
        \caption{Log in error}
        \label{enter-label}
    \end{minipage}%
    \begin{minipage}{0.35\columnwidth}
        \centering
        \includegraphics[width=\linewidth]{Images/Rloginerror2.png}
        \caption{Log in error}
        \label{enter-label}
    \end{minipage}
\end{figure}
\noindent If users submit invalid or incomplete login information, the system immediately highlights the specific issues such as missing email or password, to ensure proper input before authentication. When valid fields are entered but the credentials do not match any existing account, the backend returns an error response, and the app displays a “Login failed” message, prompting the user to try again and recheck their details.

\subsection{Main Page}
\begin{figure}[h!]
    \centering
    \includegraphics[width=0.35\columnwidth]{Images/Rhistorypage1.png}
    \caption{Main Page}
    \label{fig:enter-label}
    \end{figure}
\noindent The Main page displays all past mold detection results saved by the user. If no records are available, the page simply shows a “No history yet” message. 

\newpage
\begin{figure}[h!]
    \centering
    \begin{minipage}{0.35\columnwidth}
        \centering
        \includegraphics[width=\linewidth]{Images/Rhistorypage2.png}
        \caption{History}
        \label{enter-label}
    \end{minipage}%
    \begin{minipage}{0.35\columnwidth}
        \centering
        \includegraphics[width=\linewidth]{Images/Rhistorypagewhentap.png}
        \caption{History information}
        \label{enter-label}
    \end{minipage}
\end{figure}
\noindent Once detections exist, each entry appears with its name, detection status, and timestamp. When a user taps on any history item, the app opens the detailed Result Analysis page, allowing them to view the original image, detection outcome, and associated information exactly as it appeared during the initial analysis.

\subsection{Image Page}
\begin{figure}[h!]
    \centering
    \includegraphics[width=0.35\columnwidth]{Images/Ridentifypage.png}
    \caption{Image Page}
    \label{fig:enter-label}
    \end{figure}
\noindent When users click the Camera icon or Photo icon, the app either requests camera access or opens the gallery, and once an image is captured or selected, it proceeds to the Identify page or shows an “Access to camera denied” message if permission is blocked.

\newpage
\subsection{Identify Page}
\begin{figure}[h!]
    \centering
    \begin{minipage}{0.35\columnwidth}
        \centering
        \includegraphics[width=\linewidth]{Images/Rdetectpagemold.png}
        \caption{Mold detected}
        \label{fig:mold-detected}
    \end{minipage}%
    \begin{minipage}{0.35\columnwidth}
        \centering
        \includegraphics[width=\linewidth]{Images/Rdetectpagenomold.png}
        \caption{No mold detected}
        \label{fig:no-mold}
    \end{minipage}
\end{figure}
\noindent When users take or upload a photo, the Identify page processes the image by sending it to the backend through a secured API request. The backend loads the ONNX-based YOLO detection model, runs inference on the image, and returns a structured JSON response containing the prediction result. The Identify page then interprets this response and displays the appropriate outcome such as a mold detected alert or a safe message. At the same time, the system automatically stores the image URL, detection name, timestamp, and model output in Firestore, ensuring that every analysis is recorded and can be accessed later through the History page.

\subsection{Chatbot Page}
\begin{figure}[h!]
    \centering
    \includegraphics[width=0.35\columnwidth]{Images/Rchatbot.png}
    \caption{Chatbot page}
    \label{fig:enter-label}
    \end{figure}
\noindent The Chatbot page allows users to interact with Gom, an AI-powered assistant designed to answer questions about mold detection, prevention, and cleaning. When a user submits a message, the app sends the query to the backend, which forwards it to the OpenAI API for natural language processing. The API generates an appropriate response based on the user’s input, and the backend returns this message to the app, where it is displayed in the chat interface. This creates a real-time conversational experience that helps users understand mold-related issues and receive guidance directly within the application.

\subsection{Settings Page}
\begin{figure}[h!]
    \centering
    \begin{minipage}{0.35\columnwidth}
        \centering
        \includegraphics[width=\linewidth]{Images/Rsettings.png}
        \caption{Settings page}
        \label{fig:mold-detected}
    \end{minipage}%
    \begin{minipage}{0.35\columnwidth}
        \centering
        \includegraphics[width=\linewidth]{Images/Rsettingssuccess.png}
        \caption{Settings page}
        \label{fig:no-mold}
    \end{minipage}
\end{figure}
\noindent The Settings page allows users to view and manage their account information, including their username and password, all of which are retrieved from the backend upon login. When a user chooses to update their username or password, the app opens an edit dialog where the new value can be entered. After the user submits the change, the app sends an update request to the backend, which validates the input and writes the updated data to Firestore. Once the backend confirms a successful update, the app displays a success message and refreshes the account information on the page.

\section{DISCUSSION}

Throughout the development of GomMold, our team faced a wide range of challenges that affected the backend, AI/ML model, and frontend components equally. On the backend side, the most difficult aspect was deploying the server to the cloud. Although our original plan was to use Google Cloud because of our reliance on Firebase, the deployment process quickly became overwhelming. Build times were extremely long, and deployments repeatedly failed with unclear error messages that were buried in long, complex logs. Troubleshooting was especially difficult because GCP relies on many layers of configuration, such as Docker images, service accounts, IAM permissions, and runtime settings. Even after numerous attempts to fix the issues, each redeployment introduced new problems, making the process discouraging and time-consuming. We then shifted to Render, which simplified the interface but created new complications, especially Python version mismatches and model loading failures during the build stage. Render’s requirement to redeploy through GitHub also led to many experimental commits. Eventually, we switched to Railway, which provided faster builds and clearer error logs. However, even Railway required several adjustments, especially with library compatibility and model-loading logic. Integrating the backend with the frontend introduced another layer of work, as API responses and field names had to be aligned precisely to avoid UI errors.\\

The AI/ML component was equally challenging. Collecting a high-quality mold dataset required much more effort than we expected. Online images were limited, copyrighted, or inconsistent in quality, and the photos we took ourselves lacked variety in mold types and lighting conditions. Annotating these images was tedious because mold does not have sharp boundaries, meaning every small patch had to be carefully labeled. Even minor annotation mistakes reduced the model’s accuracy, causing confusion between mold, stains, and shadows. Initially, we planned to classify mold into different colors, but the dataset was too small and inconsistent to train a reliable color-based model. We also considered classifying mold by species or type, but this required expert knowledge and specialized datasets that we did not have access to. After training, the PyTorch model performed well locally but caused deployment errors on multiple cloud platforms. This forced us to convert the model to ONNX, which introduced a new set of challenges. The ONNX model produced different output structures, tensor shapes, and confidence values compared to the PyTorch version. We had to rewrite much of our post-processing code and re-export the model multiple times before the outputs became stable. Integrating the ONNX model with the backend required detailed testing, especially when sending results to the frontend, ensuring that color labels, messages, and analysis names appeared correctly.\\

On the frontend, the team faced both technical and communication-related challenges. Because different members worked on different devices, repository synchronization became a recurring problem. Some cloned versions were outdated, certain pages were missing after pulls, and duplicated folders caused confusion about the latest build. Flutter’s routing system also required careful attention, linking UI buttons to their correct pages and maintaining a smooth navigation flow needed multiple revisions. Ensuring consistent design across devices was another challenge, as Flutter’s layout system reacts differently depending on screen size, which led to spacing issues, misaligned components, and distorted images. The frontend also frequently broke whenever backend API outputs changed. For example, a slight change in the JSON format could cause the history page to display empty results or load incorrect images. Fixing these mismatches required constant communication and collaboration between frontend and backend members.\\

Overall, the combination of backend deployment issues, AI model preprocessing difficulties, and frontend synchronization problems made this project a complex but meaningful learning experience. Despite the challenges, the project helped us build confidence in cloud deployment, dataset preparation, machine-learning integration, version control, and cross-team communication. Since this was our first large-scale development project, much of our time was spent learning by doing, which naturally led to mistakes but also to significant growth. Moving forward, we plan to adopt clearer communication methods, establish a more structured Git workflow, and perform regular code reviews to ensure smoother integration. These improvements will help us work more efficiently in future projects and avoid many of the issues we encountered during GomMold’s development.



\begin{thebibliography}{00}

\bibitem{b1} 
App Store. (n.d.). Mould Identifier \& AI Scanner App. https://apps.apple.com/kr/app/mould-identifier-ai-scanner/id6743941888?l=en-GB

\bibitem{b2} 
App Store. (n.d.). Mold Finder AI, Inspection App. https://apps.apple.com/kr/app/mold-finder-ai-inspection/id6689505239?l=en-GB

\bibitem{b3} 
AND Academy. Stevens, E. (2025, August 14). What is Figma – uses, benefits, and features. https://www.andacademy.com/resources/blog/graphic-design/what-is-figma/

\bibitem{b4} 
Back4App Blog. Team, B. E. (2021, September 9). An overview of Firebase BaaS. https://blog.back4app.com/backend-as-a-service-firebase/

\bibitem{b5} 
Bergmann, D. (2025, October 23). What is machine learning? IBM. https://www.ibm.com/think/topics/machine-learning

\bibitem{b6} 
Bergmann, D., \& Stryker, C. (2025, October 20). What is PyTorch? IBM. https://www.ibm.com/think/topics/pytorch

\bibitem{b7} 
Build apps for any screen. Flutter. (n.d.). https://flutter.dev/

\bibitem{b8} 
Dart Programming Language: Top features and applications. Khan, S. (2025, October 17). Xavor Corporation. https://www.xavor.com/blog/dart-programming-language-features-and-applications/

\bibitem{b9} 
EBSCO Research Starters. (n.d.). Windows operating system. https://www.ebsco.com/research-starters/computer-science/windows-operating-system

\bibitem{b10} 
Expert Software Development | Kaopiz Holdings. Hồng, H. N. (2025, May 26). What is OpenAI? https://kaopiz.com/en/articles/what-is-openai/

\bibitem{b11} 
Firebase. Google. (n.d.). https://firebase.google.com/

\bibitem{b12} 
Flutter Architectural Overview. Flutter Docs. (n.d.-b). https://docs.flutter.dev/resources/architectural-overview

\bibitem{b13} 
Flask Documentation. Welcome to Flask (3.1.x). https://flask.palletsprojects.com/en/stable/

\bibitem{b14} 
Fox, L. (2023, September 17). Why should you use Flask Framework for web development? Medium. https://medium.com/@lauren-fox/why-should-you-use-flask-framework-for-web-development-f5a7233e17a6

\bibitem{b15} 
GitHub Introduction. Holcombe, J. (2025, August 28). Kinsta. https://kinsta.com/blog/what-is-github/

\bibitem{b16} 
Google Cloud Overview. Google. (n.d.-a). https://cloud.google.com/docs/overview

\bibitem{b17} 
Google. (n.d.). Add the Firebase Admin SDK to your server. https://firebase.google.com/docs/admin/setup

\bibitem{b18} 
Hot Reload. Flutter Docs. (n.d.-c). https://docs.flutter.dev/tools/hot-reload

\bibitem{b19} 
IBM. Gillis, A. S. (2025, June 11). What is macOS? WhatIs. https://www.techtarget.com/whatis/definition/Mac-OS

\bibitem{b20} 
ITarian Blog. (n.d.). What is macOS? Features, Benefits \& Latest Version. https://www.itarian.com/blog/what-is-mac-os/

\bibitem{b21} 
Kaopiz. Hồng, H. N. (2025, May 26). What is OpenAI? https://kaopiz.com/en/articles/what-is-openai/

\bibitem{b22} 
Microsoft. (2021, November 3). Visual Studio Code – Code editing, redefined. https://code.visualstudio.com/

\bibitem{b23} 
Mold Detection Dataset. Roboflow Universe. https://universe.roboflow.com/dissertationproject/mould-detection-aaron

\bibitem{b24} 
Mold Identifier: AI Scan App. Worlitzer, M. (n.d.). App Store. https://apps.apple.com/kr/app/mold-identifier-ai-scan/id6749549375?l=en-GB

\bibitem{b25} 
MoldAI App. Y. P. C. Limited. (n.d.). App Store. https://apps.apple.com/kr/app/moldai/id6689520833?l=en-GB

\bibitem{b26} 
Nguyen, A. (2025, May 5). Flask Framework: Why developers love it for web projects. Stepmedia Software. https://stepmediasoftware.com/blog/what-is-flask-used-for/

\bibitem{b27} 
Notion. (n.d.). The all-in-one workspace for your notes, tasks, wikis, and databases. https://www.notion.so/index.html

\bibitem{b28} 
Notion AI Workspace. (n.d.-a). https://www.notion.so/notion/What-s-New-157765353f2c4705bd45474e5ba8b46c

\bibitem{b29} 
Overleaf, Online LaTeX Editor. (n.d.). https://www.overleaf.com/

\bibitem{b30} 
Patel, B. (2025, May 29). Top 10 benefits of using OpenAI in web app development. Space-O Technologies. https://www.spaceotechnologies.com/blog/benefits-of-using-openai-web-app-development/

\bibitem{b31} 
Postman. (n.d.). The world’s leading API platform. https://www.postman.com/

\bibitem{b32} 
PyTorch. (n.d.). https://pytorch.org/

\bibitem{b33} 
Python.org. (n.d.). What is Python? https://www.python.org/doc/essays/blurb/

\bibitem{b34} 
Railway Docs. (n.d.). https://docs.railway.com/

\bibitem{b35} 
Roboflow Universe. Wall Computer Vision Dataset. https://universe.roboflow.com/jisupark/wall-yqazk

\bibitem{b36} 
Safi, A. (n.d.). Understanding BaaS: Firebase vs. Supabase. https://abdulkadersafi.com/blog/understanding-baas-a-deep-dive-into-firebase-vs-supabase

\bibitem{b37} 
Ultralytics YOLO Docs. (2025, November 23). https://docs.ultralytics.com/

\bibitem{b38} 
Windows. Digital Unite. https://www.digitalunite.com/technology-guides/computer-essentials/windows-10/what-windows

\end{thebibliography}



\vspace{12pt}


\end{document}